% §2.4 — NOₓ Process Dynamics
% Source: Thesis/secs/3-governing_eqns/3-NOx_proc_dyn.tex
% ===========================================================
\subsection{$\nox$ Process Dynamics \label{sec::nox_proc_dyn}}

The $\nox$ process dynamics involve only the standard SCR reaction. The gaseous $\nox$ that enters the catalyst chamber
reacts with the adsorbed ammonia on the surface forming $\nitrogen$ and $\water$ (Eley-Rideal mechanism). The volumetric rate
of $\nox$ reduction is:
\begin{align}
    \dot{\con{\nox}}^{scr} &= -k_{s2v} k_{scr} \con{\ammonia}^{ads} \con{\nox}^{in}
    \label{eqn::nox_rate}
\end{align}

% --- Molar conservation ---
\subsubsection{Molar Conservation at the Residence-Time Scale}
At the end of every residence time, a fresh parcel of gaseous reactants enters the catalyst chamber. The molar
conservation for $\nox$ reduction relates the inlet and outlet concentrations across one residence time:
\begin{align}
    \con{\nox}^{out} (k + (i+1)\tau) &= \con{\nox}^{in} (k + i\tau) + \int_{0}^{\tau}
        \dot{\con{\nox}}^{scr} (k + i\tau) \, dt
    \label{eqn::nox_molar_cons}
\end{align}
Note that the outlet $\nox$ measurement depends only on the inlet $\nox$ one residence time earlier — there is no
integrating effect of $\nox$ within the sample time. Introducing the zero-order hold approximation for the inlet $\nox$
concentration and the average surface concentration:
\begin{align*}
    \con{\nox}^{in} (k + i \tau) &\approx \con{\nox}^{in} (k) \qquad \forall \, i < n \\
    \con{\ammonia}^{ads} (k + i \tau) &\approx \sigma(k) \qquad \forall \, i < n
\end{align*}
we obtain the discrete $\nox$ process dynamics:
\begin{align}
    \con{\nox}^{out} (k + 1) &= \con{\nox}^{in} (k) \lr{1 - k_{s2v} k_{scr} \sigma(k) \tau}
    \label{eqn::NOX_process}
\end{align}
This result shows that the $\nox$ outlet concentration is a static map of the inlet concentration, modulated by the
average adsorbed ammonia $\sigma(k)$ and the residence time $\tau$.
